\begin{figure}[H]
\centering
\begin{tikzpicture}[font=\small,>=Stealth,scale=.92]
\coordinate (A) at (0,0); \coordinate (B) at (3,0); \coordinate (C) at (3,3);
\draw[very thick] (A)--(B)--(C)--cycle;
\draw (2.62,0)--(2.62,.38)--(3,.38);
\node[below] at (1.5,0) {$1$}; \node[right] at (3,1.5) {$1$};
\node[above left,rotate=45] at (1.5,1.5) {$\sqrt2$};
\node at (.55,.23) {$45^\circ$}; \node at (2.73,2.38) {$45^\circ$};

\coordinate (D) at (5.2,0); \coordinate (E) at (8.7,0); \coordinate (F) at (8.7,6.062);
\draw[very thick] (D)--(E)--(F)--cycle;
\draw (8.32,0)--(8.32,.38)--(8.7,.38);
\node[below] at (6.95,0) {$1$}; \node[right] at (8.7,3.03) {$\sqrt3$};
\node[above left,rotate=60] at (6.95,3.03) {$2$};
\node at (5.85,.25) {$60^\circ$}; \node at (8.38,5.2) {$30^\circ$};
\end{tikzpicture}
\caption{The $45^\circ$--$45^\circ$--$90^\circ$ and $30^\circ$--$60^\circ$--$90^\circ$ triangles. Their side ratios are familiar geometric targets that the analytic construction will later justify.}
\label{fig:special-triangles}
\end{figure}
